\chapter{Lebesgue外測度}

\paragraph{本章の目標}
Jordan外測度の「有限被覆」を「可算被覆」に置き換えて
Lebesgue外測度 $\lambda_d^*$ を定義し，
それが $\RR^d$ 上の外測度になることを示す．
本章ではまだ可測集合を選別しない．
$\lambda_d^*$ は全ての部分集合に定義されるが，
現時点で分かるのは可算劣加法性までである．

\section{動機づけ}

第2章では，Jordan測度は有限個の長方体による近似に基づく理論であることを見た．
そのため
\[
A = \QQ \cap [0,1]
\]
のような集合はJordan可測にならなかった．
有限個の区間で $A$ を覆おうとすると，
$A$ が $[0,1]$ に稠密であるため，
結局 $[0,1]$ 全体に近い長さを支払わされるからである．

しかし，もし可算個の区間を使ってよいなら事情は変わる．
$A=\{q_1,q_2,\dots\}$ と番号づけて，
各点 $q_n$ を長さ $\eps 2^{-n}$ の区間で覆えば，
全体の長さの和は
\[
\sum_{n=1}^{\infty} \eps 2^{-n} = \eps
\]
になる．
したがって，$\QQ \cap [0,1]$ は「長さ $0$ の集合」と見なすのが自然である．

この観察は，Jordan理論の限界が
「被覆が有限個に制限されていること」にあることを示している．
そこで，有限被覆を可算被覆に置き換えて得られる量を
Lebesgue外測度と呼ぶ．

\section{定義}

以後，次元 $d \in \NN$ を固定する．
前章と同様に，$\RR^d$ の有界長方体 $R$ の体積を $|R|$ で表す．

\begin{definition}[Lebesgue外測度]
集合 $A \subset \RR^d$ に対し，$A$ のLebesgue外測度 $\lambda_d^*(A)$ を
\[
\lambda_d^*(A)
:=
\inf\left\{
\sum_{n=1}^{\infty} |R_n|
\;\middle|\;
A \subset \bigcup_{n=1}^{\infty} R_n,\ 
R_n \text{ は有界長方体}
\right\}
\]
で定める．
\end{definition}

\begin{remark}
有限個の被覆は可算個の被覆の特別な場合である．
したがって，Lebesgue外測度はJordan外測度よりも
柔軟な近似を許している．
\end{remark}

\begin{proposition}[Jordan外測度との比較]
有界集合 $A \subset \RR^d$ に対して
\[
\lambda_d^*(A) \le m_J^*(A)
\]
が成り立つ．
\end{proposition}
\begin{proof}
$A \subset E$ を満たす基本集合 $E$ を取る．
第2章より，$E$ は有限個の互いに素な半開長方体の合併
\[
E = \bigsqcup_{j=1}^N Q_j
\]
と書け，
\[
m_J(E) = \sum_{j=1}^N |Q_j|
\]
であった．
この有限個の長方体の族は $A$ の可算被覆でもあるから，
\[
\lambda_d^*(A) \le \sum_{j=1}^N |Q_j| = m_J(E)
\]
を得る．
$E$ の取り方について下限を取ればよい．
\end{proof}

\section{基本定理}

\begin{definition}[外測度]
集合 $X$ 上の写像
\[
\mu^* : \mathcal P(X) \to [0,\infty]
\]
が外測度であるとは，次の3条件を満たすことをいう．
\begin{enumerate}
    \item $\mu^*(\emptyset)=0$
    \item $A \subset B$ ならば $\mu^*(A) \le \mu^*(B)$
    \item 任意の集合列 $\{A_n\}_{n=1}^{\infty}$ に対して
    \[
    \mu^*\left(\bigcup_{n=1}^{\infty} A_n\right)
    \le
    \sum_{n=1}^{\infty} \mu^*(A_n)
    \]
    が成り立つ
\end{enumerate}
\end{definition}

\begin{theorem}[Lebesgue外測度は外測度]
$\lambda_d^*$ は $\RR^d$ 上の外測度である．
\end{theorem}
\begin{proof}
\begin{enumerate}
    \item $\lambda_d^*(\emptyset)=0$：
    まず外測度は定義から非負なので $\lambda_d^*(\emptyset)\ge 0$ である．
    一方，任意の点 $a \in \RR^d$ に対して
    \[
    \emptyset \subset \bigcup_{n=1}^{\infty} \{a\}
    \]
    であり，$\{a\}$ は体積 $0$ の長方体だから
    \[
    \lambda_d^*(\emptyset) \le \sum_{n=1}^{\infty} 0 = 0.
    \]
    よって $\lambda_d^*(\emptyset)=0$ である．

    \item 単調性：
    $A \subset B$ とする．
    $B$ の任意の可算被覆
    \[
    B \subset \bigcup_{n=1}^{\infty} R_n
    \]
    はそのまま $A$ の可算被覆でもある．
    したがって，$A$ に対する下限は $B$ に対する下限以下であり，
    \[
    \lambda_d^*(A) \le \lambda_d^*(B)
    \]
    を得る．

    \item 可算劣加法性：
    集合列 $\{A_n\}_{n=1}^{\infty}$ を取る．
    もし右辺が $\infty$ なら不等式は自明である．
    右辺が有限であるとし，$\eps > 0$ を任意に取る．
    各 $n$ に対して，$\lambda_d^*(A_n)$ の定義より
    $A_n$ を覆う可算個の長方体 $\{R_{n,k}\}_{k=1}^{\infty}$ を
    \[
    A_n \subset \bigcup_{k=1}^{\infty} R_{n,k},
    \qquad
    \sum_{k=1}^{\infty} |R_{n,k}|
    <
    \lambda_d^*(A_n) + \frac{\eps}{2^n}
    \]
    となるように取れる．
    すると二重列 $\{R_{n,k}\}_{n,k}$ は可算個の長方体からなり，
    \[
    \bigcup_{n=1}^{\infty} A_n
    \subset
    \bigcup_{n=1}^{\infty}\bigcup_{k=1}^{\infty} R_{n,k}
    \]
    であるから，
    \begin{align*}
    \lambda_d^*\left(\bigcup_{n=1}^{\infty} A_n\right)
    &\le
    \sum_{n=1}^{\infty}\sum_{k=1}^{\infty} |R_{n,k}| \\
    &<
    \sum_{n=1}^{\infty}
    \left(
    \lambda_d^*(A_n)+\frac{\eps}{2^n}
    \right) \\
    &=
    \sum_{n=1}^{\infty}\lambda_d^*(A_n)+\eps.
    \end{align*}
    $\eps > 0$ は任意だから
    \[
    \lambda_d^*\left(\bigcup_{n=1}^{\infty} A_n\right)
    \le
    \sum_{n=1}^{\infty}\lambda_d^*(A_n)
    \]
    が従う．
\end{enumerate}
\end{proof}

\begin{corollary}[有限劣加法性]
任意の $A,B \subset \RR^d$ に対して
\[
\lambda_d^*(A \cup B) \le \lambda_d^*(A) + \lambda_d^*(B)
\]
が成り立つ．
\end{corollary}
\begin{proof}
上の定理で
\[
A_1=A,\qquad A_2=B,\qquad A_n=\emptyset\ (n \ge 3)
\]
とおけばよい．
\end{proof}

\begin{proposition}[平行移動不変性]
任意の $A \subset \RR^d$ と $c \in \RR^d$ に対して
\[
\lambda_d^*(A+c)=\lambda_d^*(A)
\]
が成り立つ．
\end{proposition}
\begin{proof}
$A \subset \bigcup_{n=1}^{\infty} R_n$ ならば
\[
A+c \subset \bigcup_{n=1}^{\infty} (R_n+c)
\]
であり，平行移動は各辺の長さを変えないので
\[
|R_n+c| = |R_n|
\]
である．
したがって
\[
\lambda_d^*(A+c) \le \sum_{n=1}^{\infty} |R_n|
\]
となるから，$A$ の被覆について下限を取って
\[
\lambda_d^*(A+c) \le \lambda_d^*(A)
\]
を得る．
逆向きの不等式は $-c$ だけ平行移動すれば従う．
\end{proof}

\section{具体例}

\begin{proposition}[長方体による上からの評価]
$A \subset R$ を満たす有界長方体 $R$ に対して
\[
\lambda_d^*(A) \le |R|
\]
が成り立つ．
特に，有界集合のLebesgue外測度は有限である．
\end{proposition}
\begin{proof}
$R$ 自身による1個の被覆を考えればよい．
\end{proof}

\begin{proposition}[1点集合は零集合]
任意の $a \in \RR^d$ に対して
\[
\lambda_d^*(\{a\})=0
\]
が成り立つ．
\end{proposition}
\begin{proof}
非負性より $\lambda_d^*(\{a\}) \ge 0$ である．
任意の $\eps > 0$ に対し，$\delta > 0$ を十分小さく取って
\[
(2\delta)^d < \eps
\]
とする．
すると
\[
\{a\}
\subset
\prod_{j=1}^d (a_j-\delta,a_j+\delta)
\]
であり，右辺の体積は $(2\delta)^d$ だから
\[
\lambda_d^*(\{a\}) \le (2\delta)^d < \eps.
\]
$\eps$ は任意なので $\lambda_d^*(\{a\})=0$ である．
\end{proof}

\begin{theorem}[可算集合は零集合]
可算集合 $A \subset \RR^d$ に対して
\[
\lambda_d^*(A)=0
\]
が成り立つ．
\end{theorem}
\begin{proof}
$A=\{a_1,a_2,\dots\}$ と書く．
上の命題と可算劣加法性より
\[
\lambda_d^*(A)
\le
\sum_{n=1}^{\infty} \lambda_d^*(\{a_n\})
=
\sum_{n=1}^{\infty} 0
=0.
\]
非負性と合わせて $\lambda_d^*(A)=0$ を得る．
\end{proof}

\begin{example}
\begin{enumerate}
    \item $\QQ^d$ は可算だから $\lambda_d^*(\QQ^d)=0$
    \item 特に
    \[
    \lambda_1^*(\QQ \cap [0,1]) = 0
    \]
    である
    \item $\ZZ^d$ も可算だから $\lambda_d^*(\ZZ^d)=0$
\end{enumerate}
\end{example}

\begin{remark}
第2章でJordan可測でなかった
\[
\QQ \cap [0,1]
\]
が，Lebesgue外測度では長さ $0$ になる．
ここに「有限被覆」と「可算被覆」の差が最もはっきり現れる．
\end{remark}

\begin{remark}
一方で，この段階ではまだ
\[
\lambda_d^*(R)=|R|
\]
を一般の長方体 $R$ について証明したわけではない．
今示したのは，上からの評価
\[
\lambda_d^*(R) \le |R|
\]
だけである．
逆向きの評価や，加法性の成立は，
次章以降で内測度や可測性を導入して初めて明確になる．
\end{remark}

\section{解釈}

Lebesgue外測度の本質は，
Jordan外測度の有限被覆を可算被覆へ緩めたことにある．
この1つの変更によって，
\begin{enumerate}
    \item 可算集合のような「薄い集合」を体積 $0$ として扱える
    \item 有界性を仮定せず，全ての部分集合に量を割り当てられる
    \item しかしその代わり，現時点では加法性は失われ，劣加法性しか残らない
\end{enumerate}
という新しい状況が生じる．

したがって $\lambda_d^*$ は，まだ「測度」ではなく「外測度」である．
次章ではこれに対応する内測度を導入し，
\[
\text{外測度}=\text{内測度}
\]
が成り立つ集合をLebesgue可測集合として定義する．
